\chapter*{What Has Been Built}

This work set out to answer a single question: what is the minimum
mathematical structure needed to ground physics, knowledge, life, and
mind?
The answer, developed across six parts, is that a graph-structured lambda
theory --- three ingredients: a grammar, equations, and rewrite rules ---
suffices, provided it names the site at which its terms interact and puts
a meter on that site.

Part~\ref{part:machinery} built the apparatus.
Two strengthenings of a very thin starting notion give an interactive
theory and then a continued interactive one, and on the category of the
latter sit two constructions used everywhere afterwards: a cost monad
that installs a meter and a history monad that installs a log.
The metered theory was exhibited concretely, as a rho calculus in which
token stacks are terms, purses are located, and a computation is accepted
in advance on a linear proof that it can pay.
The virtual token was shown to exist exactly when no cycle of conversions
pays for itself.
And the logic a theory generates was given as far as it has been built,
with an honest note saying which of its ingredients are still described
rather than exhibited.

Part~\ref{part:ecology} built the knower.
A learner is a term among terms with a bounded supply of tokens and no
way to mint more, and two things follow that were not designed in: that a
specimen affords one measurement, so hypotheses must be about kinds; and
that being right is not the same as being fed.
The unit of learning then had to move four times --- individual, lineage,
composite, network --- because an individual revising its theory in small
steps cannot leave the ball it started in.
At the top the picture closed on itself: the network is an individual,
and an individual can be a scientist.
This is the most explicitly worked material in the book and the only
place where a claim is checked against a number.

Part~\ref{part:physics} asked what world such a learner finds itself in,
and answered that it is already in one, and that it is one of many.
Each axiom imposed on a cost-accounted theory names a subcategory:
reversibility gives time symmetry, recorded histories give a causal order
and a proper time, a decoration into a semiring gives nondeterminism or
cost or probability or amplitude, and arbitrage-free conversion gives a
conserved quantity.
When conversion is lossy the conserved quantity becomes a free energy and
the second law arrives on schedule, and the erasure a computation must
perform turns out to be bounded below by the same holonomy that measures
the failure of conservation.
These are constructions and theorems, not analogies --- and where what
remains is a conjecture, chiefly the identification of the conserved
charge with a generator of time evolution, we have said so.
The ladder descends some distance towards our own physics and stops at an
obstruction: complex amplitudes give interference and do not by
themselves give the Born rule.

Part~\ref{part:sciencebot} asked how much of that world gets found, and
established that agents embedded in a massive population will, under any
finite experimental budget, converge to a theory of the cluster hierarchy
rather than the bisimulation quotient.
This is not a failure of the scientific method.
The cluster hierarchy is objectively present, arising from the
communication geometry.
A budget-limited agent does good science and arrives at correct theories
at its scale, unaware that the scale is not the bottom.
The compactness argument --- that compact populations support consensus,
and that coherence clusters are maximal compact subpopulations --- is new
to our knowledge and is the sharpest result of that part.

Part~\ref{part:choice} asked what a mind can do that a Turing machine
cannot, and located the answer at the apertures.
Metering subtracts power, logging adds none, and imposing axioms carves
out subcategories, so nothing in the first four parts leaves the Turing
degree it started in.
What does leave it is the resolution of races: picking a winner is a
choice function, a complete schedule is an instance of the axiom of
choice, and the strength required is graded.
The oracle was never something to bolt on; in any calculus where
interaction is primitive and confluence fails, it is present at every
cut, wearing the name \emph{scheduler}.
On that footing sentience was defined as a graded feeling state within a
stage and consciousness as access to a higher one --- the first
continuous, the second not, and the two linked by cycles that can run in
either direction.

Part~\ref{part:origins} grounded the origins of life.
Assembly index is depth in the open synchronization tree.
Copy number is count within a namespace, taken at the resolution the
counter can afford.
Life emerges when a prebiotic population wanders into the basin of
attraction of a replicating fixed point, which is dense in the logical
topology, and the phase transition is the crossing of that boundary.
The part closed by introducing two constructions that had not previously
met.
Each level of the medium tower of Part~\ref{part:ecology} costs at least
one joining step, so the tower is bounded by the assembly index; the
depth of a learner's affordable hypotheses is bounded by the tower; and
modal nesting depth characterises $n$-step bisimulation.
An ecology-as-mind assembled by $n$ joining steps can therefore resolve
its world at best to $n$-step bisimulation.
Mind is bounded by manufacture.
And read backwards the same bound says what a mind costs: to resolve at a
depth is to have been built to it, to leave the basin one started in is to
carry a germ and pay for a level of the tower, and to use a deep
architecture at all is to hold copies enough to sustain it against the
failure rate of one's own interior.

None of these correspondences are metaphors.
Each is a mathematical construction, stated precisely enough to be wrong.
